\documentclass[12pt]{article}

% --- PAQUETES BÁSICOS Y IDIOMA ---
\usepackage[utf8]{inputenc}
\usepackage[spanish]{babel}

% --- TIPOGRAFÍA COURIER (Tipo máquina de escribir de la imagen) ---
\usepackage{courier}
\renewcommand{\familydefault}{\ttdefault} % Configura Courier como fuente principal

% --- MÁRGENES MÁS GRANDES ---
\usepackage{geometry}
\geometry{
    letterpaper,
    left=1.5in,   % Margen izquierdo amplio para guion
    right=1.5in,  % Margen derecho amplio
    top=1.2in,    % Margen superior para dar aire al encabezado
    bottom=1.2in  % Margen inferior para el pie de página
}

\usepackage{ragged2e}

% --- CONFIGURACIÓN DE ENCABEZADOS Y PIES DE PÁGINA (Estilo image.png) ---
\usepackage{fancyhdr}
\usepackage{datetime} % Para poder manipular la hora si fuera necesario

\pagestyle{fancy}
\fancyhf{} % Borra los estilos por defecto

% Grosor de las líneas (las dejamos en 0pt para que no haya líneas divisorias, igual que en la foto)
\renewcommand{\headrulewidth}{0pt}
\renewcommand{\footrulewidth}{0pt}

% Fuentes más chicas para los encabezados y pies (\scriptsize o \footnotesize)
\newcommand{\estiloHeader}[1]{\textsf{\scriptsize #1}}

% Parte Superior (Header)
\lhead{\estiloHeader{16/10/23, 21:48}}             % Superior Izquierda: Fecha y Hora fija de la foto (o puedes usar \today)
\chead{\estiloHeader{Celtx - Script My Project}}    % Superior Centro: Nombre del archivo / Proyecto

% Parte Inferior (Footer)
\lfoot{\estiloHeader{about:blank}}                  % Inferior Izquierda: Título/Ruta (en tu foto dice about:blank)
\rfoot{\estiloHeader{\thepage/\pageref{LastPage}}}  % Inferior Derecha: Número de página formato X/Y

% Paquete necesario para calcular el total de páginas (el "/10" de la imagen)
\usepackage{lastpage}


% --- CONFIGURACIÓN DE FORMATO DE GUION ---
% (Ajustado para que mantenga el formato Courier y la alineación correcta)
\newenvironment{scenedesc}{\vspace{1em}\par\noindent}{\vspace{1em}}

\newenvironment{dialogue}[1]{
    \vspace{0.5em}\par
    \begin{center}\textbf{\MakeUppercase{#1}}\end{center}
    \vspace{-0.5em}
    \begin{list}{}{
        \leftmargin=1.0in
        \rightmargin=1.0in
        \parsep=0pt
    }\item[]
}{
    \end{list}\vspace{0.5em}
}
\newcommand{\parenthetical}[1]{\vspace{-0.3em}\begin{center}(#1)\end{center}\vspace{-0.3em}}


% --- INICIO DEL DOCUMENTO ---
\begin{document}
\raggedright % Mantiene el alineado a la izquierda sin justificar, idéntico a los guiones

% Nota: Si usas portadas independientes, heredarán este diseño. 
% Si no quieres encabezados en la portada, puedes usar \thispagestyle{plain} dentro de ella.
\begin{titlepage}
    
\RaggedRight

\thispagestyle{empty}
\vspace*{2.5in}
\begin{center}
    {\Huge \textbf{[Ciudad Vicio]}}\\[1cm]
    {\large Borrador: SINOPSIS, ESCALETA Y PRIMER ESCENA}\\[2cm]
    Escrito por\\
    \textbf{[CINCOSEIS]}
\end{center}

\vfill
\noindent SERGIO GONZÁLEZ\\
Contacto: [ suglugr@gmail.com]\hfill \today
\newpage
\newpage


\end{titlepage}

% --- INICIO DE LA ESCALETA ---
\pagestyle{plain}
\setcounter{page}{1}

% Configuramos que los saltos de párrafo dejen una línea en blanco automática
\setlength{\parskip}{\baselineskip} 

\textbf{SINOPSIS}

Tema: VENGANZA

LOGLINE: 
Amaury es un joven que junto a su mentor Teódulo realiza pequeñas estafas.
Juntos roban dinero a un hombre "M.G.", un poderoso político corrupto, quién asesina a Teódulo en represalia por el robo. Antes de morir, Teódulo le dice a Amaury que busque a un estafador legendario llamado El Peluka. Peluka y Amaury se unen para vengar la muerte de Teódulo estafando a "M.G.".

\textbf{SINOPSIS LARGA}

Amaury está enviciado con el skateboarding, su estética, su mentalidad y su libertinaje. Convencido de que los atributos que lo hacen buen patinador son un camino fácil al éxito se vuelve un hábil estafador junto a su mentor Tédulo, un viejo que conoció patinando. 

Juntos hacen una estafa a uno de los hombres de M.G., un político corrupto dueño de casas de cambio famosas de lavar dinero; Lo engañan y le quitan un sobre de dinero, se dividen las ganancias y revelan que trabajan juntos.

Amaury recibe su parte, que es hasta ahora es la mayor cantidad de dinero que ha ganado y piensa que es buena idea salir de fiesta, entre apuestas y alcohol gasta todo el dinero esa misma noche. 


al regresar triste y enojado, Teódulo, su mentor, le explica su retiro y le cuenta de Peluka, un estafador de leyenda que hace parecer su paga de miles como centavos. 

Las acciones de Amaury la noche anterior llaman la atención de la calle. M.G. al enterarse manda asesinar a Teódulo, quién asociaron como el autor de la estafa. Por otro lado, contrata a Sáncho, un policía, para identificar  y amenazar al cómplice(Amaury). 

M.G. desconoce a Amaury, sin embargo, Sáncho sabe bien quien es. Tomando ventaja de esto, Sáncho le exige a Amaury una parte del dinero robado bajo amenaza de arresto. Después de ser robado y golpeado por Sáncho, Amaury se entera del asesinato del viejo Teódulo y decide huir en búsqueda de Peluka.

Amaury conoce a Eliseo Peluka, un patinador mayor que ha perdido jovialidad y parece ser irresponsable y libertino. Tras una noche de fiesta, Amaury encuentra a Peluka crudo y lleno de botellas. Enojado por su actitud, Amaury lo convence de ser su mentor y su participación en una estafa elaborada. 

Peluka, sentido por el remordimiento del duelo de la muerte de su amigo accede a mentorear a Amaury. Juntos planean una venganza en contra de M.G.

Después de un tiempo de investigación descubren su adicción a las apuestas y su afición por el skateboard. M.G. mantiene un perfil bajo en el mundo skate, se relaciona exclusivamente en spots o skateparks. Peluka y Amaury reúnen a un equipo y planean engañarlo montando un club de apuestas falso. 

El primer contacto se da cuándo se infiltran en uno de sus juegos de Póker. El famoso juego de Póker en dónde se mueven miles de pesos y nunca pierde el anfitrión M.G. El secreto es que M.G. soborna al repartidor de cartas para usar una baraja marcada y solo juega 3 rondas. 


`Entrada de 70 mil pesos y un juego de skate` Son las dos reglas que revela la investigación del dúo. Peluka descubre que hay un código de comunicación entre los participantes con trucos del skate. Cada combinación de cartas lleva el nombre de un truco skate. 

Durante el juego de cartas, pasa lo previsto, el repartidor de cartas esta sobornado por M.G. y controla la cartas, pero Peluka cambia sus cartas y gana el juego. M.G. esta furioso por que sabe que perdió ante una trampa pero delatar a Peluka lo delataria a el también. Amaury es enviado a cobrar el dinero de la apuesta al camerino de M.G. 

Cuando llega por el dinero Amaury le explica a M.G. como Peluka hizo trampa lo que hace estar furioso, endeudado y humillado por culpa de Peluka. Amaury lo convence de estar de su lado y de su enemistad con Peluka. Le dice que es mesero en su club de apuestas y que ha estado engañando a Peluka por tiempo robándole en las apuestas con un sistema. Él siempre es el ganador de los eventos transmitidos por PXTV, pues al tener un retardo en la señal, él puede saber 3 minutos antes del resultado de las peleas, juegos o deportes. Como tiene esta información antes la apuesta siempre resultará ganadora.

Sáncho le sigue el rastro a Amaury y lo encara, pero es intervenido por Levi, un agente encubierto de seguridad nacional. Levi le explica a Sáncho que llevan mucho tiempo buscando a Peluka y que Amaury es la persona que los llevará a arrestar a Peluka, con la finalidad de que Peluka no sepa nada, esto lo están haciendo "bajo el agua". Gracias a Sáncho, Amaury es interceptado por Levi quien lo obliga a traicionar a Peluka.

Al final cuando M.G. da 12 millones de pesos en efectivo y van a dar el resultado de la apuesta entran policías en cubierto junto a Sáncho, al ver esto Amaury y Peluka se disparan entre ellos, Sáncho ve a M.G. y lo saca del lío pues no puede ser visto ahí y se revela que los policías encubiertos son falsos y que Peluka y Amaury fingieron su muerte.



% Restauramos el parskip a 0 para que la lista de la escaleta no se deforme
\setlength{\parskip}{0pt}
 
\pagebreak
\section{ Escaleta}
\begin{enumerate}
    \item Amaury y Teódulo estafan sin saberlo a un mandadero del poderoso hombre M.G. en la calle.
    \item El corrupto policía Sáncho  intercepta a Amaury y le exige una parte del dinero robado bajo amenaza de arresto.
    \item Los hombres de M.G. asesinan a Teódulo en represalia por el robo del dinero, obligando a Amaury  a huir por su vida.
    \item Amaury  viaja a Texcoco para buscar a El Peluka (Eliseo), un veterano maestro del engaño y amigo de Teódulo, para planear una venganza en contra de M.G.
    \item Eliseo acepta el caso y reúne a un talentoso y selecto grupo de estafadores profesionales para diseñar un elaborado plan contra M.G.
    \item El equipo decide montar un falso club de apuestas para engañar a M.G.
    \item Peluka se infiltra en su exclusiva partida de póker privada.
    \item Peluka vence a M.G. en el juego utilizando sus propias trampas y lo humilla frente a los demás pasajeros, provocando su furia.
    \item Amaury, usando el alias de Brayan, se acerca a M.G. fingiendo odiar a Peluka y le propone una alianza para arruinarlo.
    \item Amaury se gana la confianza de M.G. y lo convence que tiene información que hace ganar siempre en la casa de apuestas.
    \item El teniente Sáncho sigue el rastro de Amaury  en Texcoco, pero es retenido e interrogado por agentes policiacos comandados por Levi.
    \item M.G.  visita el club de apuestas falso y ganan una gran suma de dinero comprobando el método de Amaury. 
    \item Un asesino profesional contratado por M.G. persigue en secreto a Amaury por toda la ciudad para eliminarlo.
    \item El falso informante  le entrega a M.G. la supuesta jugada ganadora definitiva para la carrera final.
    \item M.G. apuesta una fortuna millonaria en la ventanilla del club falso justo antes de que comience asegurando ganar la apuesta. 
    \item El informante llega a la casa de apuestas, le dice a M.G. que hizo la apuesta de forma errónea.
    \item M.G. entra en pánico e intenta recuperar desesperadamente su dinero pero aparecen los federales  y Sáncho en el club.
    \item Levi obliga a Amaury a traicionar a Peluka, lo que provoca un tiroteo fingido donde ambos estafadores simulan caer muertos.
    \item Levi se lleva a M.G. a toda prisa de la escena del crimen para evitar que lo vinculen con el escándalo y los supuestos asesinatos.
    \item Amaury y Peluka se levantan ilesos, celebran el éxito absoluto de la gran estafa y desmantelan rápidamente el club falso.
\end{enumerate}

\begin{itemize}
    \item 
\end{itemize}

\newpage

\begin{center}
    \LARGE \textbf{CIUDAD VICIO}  \\
    \vspace{1em}
    \small Primer Escena: Estafa cambio de pañuelo
\end{center}
\raggedright

\vspace{2em}

\textbf{EXT. CALLEJÓN - DÍA}

\begin{scenedesc}
Miguel sale de una puerta pequeña al fondo de un callejón  y gira a la izquierda adentrándose en una vía más amplia. Aparentemente, la calle está vacía, cuando se escuchan unas ruedas y aparece un joven sobre una patineta en dirección contraria a la de Miguel. El joven se detiene justo a su lado, y ambos se miran durante un instante. 

De repente, Miguel oye gritos procedentes de algún lugar a sus espaldas. Un hombre de mediana edad y complexión menuda dobla la esquina a toda velocidad, corriendo desesperadamente. Lleva una cartera en la mano. A lo lejos se escuchan unos gritos.
\end{scenedesc}

\begin{dialogue}{Desconocido (Teódulo)}
¡Cabrón! No te quieras pasar de verga. Esa es mi Mochila. 
\end{dialogue}

\begin{scenedesc}
Miguel se queda inmóvil, sin el menor interés en ejercer la justicia. Justo cuando el ladrón está a punto de pasar de largo, Amaury alza su tabla y se la lanza a las piernas, haciéndolo tropezar y rodar por el suelo, provocando que suelte la Mochila que llevaba.

Amaury corre y patea la mochila en la dirección de Miguel, casi por reflejo, Miguel la recoge. El ladrón se pone en pie y saca una navaja, Miguel y Amaury se alertan en posición defensiva.

El ladrón, al darse cuenta que buscan la pelea lanza un navajaso hacia Amaury pero Amaury lo esquiva con facilidad. Después de esto, el ladrón decide huir.  
\end{scenedesc}

\begin{scenedesc}
El ladrón huye hacia la calle principal y desaparece de la vista. De entre unos vochecitos viejos y abandonados entre cristales rotos y metales puntiagudos, un viejo se arrastra hacia el medio de la calle. 

Miguel al verlo que está sangrando de la pierna lo toma de las manos y le ayuda a sentarse en el suelo mientras sostiene la mochila en el hombro. Amaury se acerca al viejo con indiferencia a Miguel.  
\end{scenedesc}
\begin{dialogue}{Miguel}
\parenthetical{Sosteniendo al viejo}
¿Está bien? Déjeme darle una ayudadita Don.
\end{dialogue}

\begin{dialogue}{Hombre Viejo (Teódulo)}
 Mi mochila y mi dinero.
\end{dialogue}

\begin{dialogue}{Miguel}
\parenthetical{Entregándosela}
Tranquilo, Don. Aquí la tengo
\end{dialogue}

\begin{dialogue}{Hombre Viejo (Teódulo)}
\parenthetical{Aliviado, tomándola}
Ese pendejo me empujo contra este carro abandonado siento como que me rompí la pierna.
\end{dialogue}

\begin{dialogue}{Desconocido (Amaury)}
No se mueva jefe, mejor le hablamos a la ambulancia o a la patrulla.
\end{dialogue}

\begin{dialogue}{Hombre Viejo (Teódulo)}
¡No, la tira  no!
\end{dialogue}

\begin{scenedesc}
Teódulo abre la mochila, revelando un fajo de billetes 500 pesos. Miguel se emociona al ver el dinero fácil. Teódulo se intenta poner de pie.
\end{scenedesc}

\begin{dialogue}{Hombre Viejo (Teódulo)}
    Me tengo que ir ya.
\end{dialogue}

\begin{dialogue}{Miguel}
    ¿Cómo? No se puede ir asi jefe.  
\end{dialogue}


\begin{dialogue}{Hombre Viejo (Teódulo)}

    No lo entiendes, es ahuevo que tengo que llegar. Yo soy el encargado del puesto de aqui la esquina, esa feria es de la unión y es la segunda vez que no llevo el corte de la semana a tiempo. Hoy me mandaron un mensajito con un tipo en moto. Si no entrego este dinero en el punto antes de las 6 van a decir que me pelé con todo y feria. Me van a quebrar.

\end{dialogue}


\begin{dialogue}{Desconocido (Amaury)}
¡Viejo!, no puede ni caminar. Ni tirándole paro llegamos a tiempo.
\end{dialogue}

\begin{dialogue}{Hombre Viejo (Teódulo)}
Tengo 3 mil si me entregan el dinero. Se van a michas.
\end{dialogue}

\begin{dialogue}{Desconocido (Amaury)}
\parenthetical{vacila, visiblemente nervioso}
¡Saco! Yo creo ya me campanearon los de la cuadra. ¿Que tal que me estan esperando?  y traigo la bronca.  
\end{dialogue}

\begin{dialogue}{Hombre Viejo (Teódulo)}
Nadie va a saber que llevas.
\end{dialogue}

\begin{dialogue}{Desconocido (Amaury)}
\parenthetical{sacudiendo la cabeza}
Si quiere le ayudo a acomodarse, pero le saco al encarguito. 
\end{dialogue}

\begin{dialogue}{Hombre Viejo (Teódulo)}
\parenthetical{desesperado, girándose hacia Miguel}
¿Que me dices tu cabrón? Te llevas los 3 completos pero hazme el favor.
\end{dialogue}

\begin{dialogue}{Desconocido (Amaury)}
No manche Don, ¿por que va a confiar en él? Esta caliente esa bronca y es mucho varo. 
\end{dialogue}

\begin{dialogue}{Miguel}
Chamaco miado, a ti que te valga verga. 
\parenthetical{se dirije a Teódulo}
¿Cual punto es?
\end{dialogue}

\begin{dialogue}{Hombre Viejo (Teódulo)}
Es junto a la capilla del centro, la casa azul, tocas en el zaguán negro. Hay veinte mil ahí, y aquí tres mil para usted.
\end{dialogue}

\begin{dialogue}{Miguel}
\parenthetical{tomando el fajo de billetes y los billetes sueltos}
Pues ya esta Don, yo me encargo, usted tranquilo. 
\end{dialogue}

\begin{scenedesc}
Miguel se guarda el dinero con total naturalidad en su chamarra junto a otro sobre con dinero proveniente de la casa de cambio. 
Amaury, que había sacado un pañuelo de su bolsillo para empezar a ``vendar'' la pierna de Teódulo, se detiene y lo observa fijamente.
\end{scenedesc}

\begin{dialogue}{Desconocido (Amaury)}
En cuanto te campaneen en la esquina va a salir el clavo. 
\end{dialogue}

\begin{dialogue}{Hombre Viejo (Teódulo)}
\parenthetical{asustado de nuevo repentinamente}
¿Como crees que lo vas a entusar asi?
\end{dialogue}

\begin{dialogue}{Desconocido (Amaury)}
Aunque sea en su calcetín maestro, ¿No tiene nada para disimular?
\end{dialogue}

\begin{dialogue}{Miguel}
No
\end{dialogue}

\begin{dialogue}{Desconocido (Amaury)}
¿Algo como un trapo, un pañuelo?
\end{dialogue}

\begin{dialogue}{Hombre Viejo (Teódulo)}
A ver aqui
\parenthetical{Teódulo saca un pañuelo blanco arrugado}
\end{dialogue}

\begin{dialogue}{Desconocido (Amaury)}
Póngame aqui el dinero
\end{dialogue}

\begin{scenedesc}
Miguel vacila un momento, pero luego de unos segundos le da el fajo de billetes a Amaury. Amaury los coloca en el centro del pañuelo.
\end{scenedesc}

\begin{dialogue}{Desconocido (Amaury)}
    Ponle todo lo que traigas man, el clavo es el clavo.
\end{dialogue}

\begin{dialogue}{Hombre Viejo (Teódulo)}
\parenthetical{suplicando}
Mi pierna, trae un doctor, tu métele prisa.
\end{dialogue}

\begin{scenedesc}
Miguel, arrastrado por la urgencia del momento, saca el dinero que traía de la casa de cambio y lo deposita también dentro del pañuelo. Amaury junta rápidamente las esquinas del pañuelo y las ata con un nudo muy firme.
\end{scenedesc}

\begin{dialogue}{Desconocido (Amaury)}
\parenthetical{haciendo una demostración al deslizar el fajo dentro de sus pantalónes}
El clavo bien clavado. Ese entuse no sale y no se ve.
\end{dialogue}

\begin{scenedesc}
Amaury saca el fajo de sus pantalones , cambiando el pañuelo falso y el real,  y le entrega el paquete falso idéntico a Miguel.
\end{scenedesc}

\begin{dialogue}{Miguel}
\parenthetical{metiéndose el fajo falso en los pantalones}
Ahuevo, si se la sabe chamaco mión.
\end{dialogue}

\begin{dialogue}{Desconocido (Amaury)}
Es sin preguntar. 
\end{dialogue}

\begin{scenedesc}
Miguel echa un vistazo a las salidas del callejón, da la vuelta y se aleja caminando lo más rápido posible sin llegar a correr. \linebreak

Amaury lo observa marcharse hasta que Miguel da vuelta en la esquina y se sube a un Taxi. En el preciso instante en que lo pierde de vista, Teódulo se incorpora por completo, olvidándose instantáneamente de su falsa cojera. Amaury saca el pañuelo real que escondía tras su espalda, lo desata y sonríe ante el gigantesco botín. Ambos comparten una breve y silenciosa carcajada antes de marcharse por caminos separados. \linebreak

Miguel, por otra parte, cuando el taxi se ha alejado del lugar, confiado de su robo saca el pañuelo de su pantalón y descubre billetes hechos de recortes de papel periódico. \\
\end{scenedesc}



\end{document}